\sectionheading{5.8 Oracle}

Oracle মূলত একটি বহুল ব্যবহৃত
\textbf{Relational Database Management System (RDBMS)}
এবং Database Platform। এর সাহায্যে বিপুল পরিমাণ Data
সুশৃঙ্খলভাবে সংরক্ষণ, অনুসন্ধান, পরিবর্তন ও পরিচালনা করা যায়।

Java বা Python-এর মতো Oracle নিজে কোনো General Purpose
Programming Language নয়। Oracle Database-এর সঙ্গে Data নিয়ে
কাজ করার জন্য প্রধানত \textbf{SQL} এবং Programming সুবিধার জন্য
\textbf{PL/SQL} ব্যবহার করা হয়।

%================================================
% Definition
%================================================
\booksubheading{Oracle কী?}

\begin{definitionbox}

\textbf{\color{darkTeal}সংজ্ঞা}

Oracle হলো একটি
\textbf{Relational Database Management System (RDBMS)},
যার সাহায্যে Database তৈরি, Data সংরক্ষণ, অনুসন্ধান, পরিবর্তন
ও পরিচালনা করা যায়।

\end{definitionbox}

\textbf{সহজভাবে মনে রাখবে—}

\begin{center}

\textbf{
Oracle
\(\;=\;\)
Database Management System
}

\end{center}

Oracle-এর প্রধান কাজ হলো Data-কে Database-এর মধ্যে
সুশৃঙ্খলভাবে সংরক্ষণ ও পরিচালনা করা।

%================================================
% RDBMS
%================================================
\booksubheading{RDBMS কী?}

\begin{definitionbox}

\textbf{\color{darkTeal}RDBMS}

যে Database Management System-এ Data সাধারণত
\textbf{Table-এর Row ও Column} আকারে সংরক্ষণ করা হয় এবং
বিভিন্ন Table-এর মধ্যে সম্পর্ক স্থাপন করা যায়, তাকে
\textbf{Relational Database Management System (RDBMS)} বলে।

\end{definitionbox}

Oracle একটি পরিচিত RDBMS।

\begin{examplebox}{সহজ উদাহরণ}

একটি Student Table হতে পারে—

\begin{center}

\begin{tabularx}{0.72\linewidth}{|Y|Y|Y|}
\hline

\rowcolor{softBlue}
\multicolumn{1}{|c|}{\textbf{Roll}}
&
\multicolumn{1}{c|}{\textbf{Name}}
&
\multicolumn{1}{c|}{\textbf{GPA}}
\\
\hline

101 & Rahim & 4.75\\
\hline

102 & Karim & 5.00\\
\hline

\end{tabularx}

\end{center}

এখানে প্রতিটি Row একটি Record এবং প্রতিটি Column একটি
নির্দিষ্ট Field বা Attribute-এর Data প্রকাশ করে।

\end{examplebox}

%================================================
% SQL
%================================================
\booksubheading{Oracle-এ SQL-এর ব্যবহার}

Oracle Database-এ Data নিয়ে কাজ করার জন্য
\textbf{SQL} বহুল ব্যবহৃত হয়।

\begin{definitionbox}

\textbf{\color{darkTeal}SQL}

SQL-এর পূর্ণরূপ হলো
\textbf{Structured Query Language}।

Database-এর Data সংরক্ষণ, অনুসন্ধান, পরিবর্তন ও পরিচালনার জন্য
SQL ব্যবহার করা হয়।

\end{definitionbox}

SQL-এর সাহায্যে—

\begin{itemize}

\item নতুন Data যোগ করা যায়।

\item প্রয়োজনীয় Data খুঁজে বের করা যায়।

\item Data পরিবর্তন করা যায়।

\item Data মুছে ফেলা যায়।

\item Table তৈরি ও পরিবর্তন করা যায়।

\end{itemize}

%================================================
% SQL Example
%================================================
\begin{examplebox}{SQL-এর ছোট উদাহরণ}

ধরা যাক, \textenglish{\ttfamily student} নামের Table-এর
সব Data দেখতে হবে।

\begin{codebox}

\begin{center}
{\textenglish{\ttfamily SELECT * FROM student;}}
\end{center}

\end{codebox}

এখানে \textenglish{\ttfamily SELECT} ব্যবহার করে Database থেকে
Data নেওয়া হচ্ছে।

\end{examplebox}

%================================================
% PL SQL
%================================================
\booksubheading{PL/SQL কী?}

\begin{definitionbox}

\textbf{\color{darkTeal}PL/SQL}

PL/SQL-এর পূর্ণরূপ হলো
\textbf{Procedural Language/SQL}।

এটি Oracle-এর SQL-এর সঙ্গে Procedural Programming-এর বিভিন্ন
সুবিধা যুক্ত করে।

\end{definitionbox}

PL/SQL-এ SQL Statement-এর পাশাপাশি—

\begin{itemize}

\item Variable

\item Condition

\item Loop

\item Procedure

\item Function

\end{itemize}

ইত্যাদি ব্যবহার করা যায়।

\textbf{সহজভাবে মনে রাখবে—}

\begin{center}

\textbf{
SQL
\(\;+\;\)
Procedural Programming
\(\;=\;\)
PL/SQL
}

\end{center}

%================================================
% SQL vs PL SQL
%================================================
\booksubheading{SQL ও PL/SQL-এর পার্থক্য}

\begin{tabularx}{\linewidth}{|>{\bfseries}p{0.22\linewidth}|Y|Y|}
\hline

\rowcolor{softBlue}
\multicolumn{1}{|c|}{\textbf{বিষয়}}
&
\multicolumn{1}{c|}{\textbf{SQL}}
&
\multicolumn{1}{c|}{\textbf{PL/SQL}}
\\
\hline

মূল উদ্দেশ্য
&
Database-এর Data নিয়ে কাজ
&
SQL-এর সঙ্গে Programming Logic ব্যবহার
\\
\hline

Statement
&
Query ও Database Operation
&
একাধিক Statement Block আকারে ব্যবহার করা যায়
\\
\hline

Variable
&
সাধারণ Query-তে Programming Variable ব্যবহারের ধারণা নেই
&
Variable ব্যবহার করা যায়
\\
\hline

Condition
&
সাধারণ SQL-এর মূল Programming বৈশিষ্ট্য নয়
&
Condition ব্যবহার করা যায়
\\
\hline

Loop
&
সাধারণ SQL-এর মূল Programming বৈশিষ্ট্য নয়
&
Loop ব্যবহার করা যায়
\\
\hline

ব্যবহার
&
Data Insert, Search, Update, Delete
&
Procedure, Function ও Logic-ভিত্তিক Database কাজ
\\
\hline

\end{tabularx}

%================================================
% Features
%================================================
\booksubheading{Oracle-এর প্রধান বৈশিষ্ট্য}

\begin{compacttwocolpoints}

\item বিপুল পরিমাণ Data সংরক্ষণ করা যায়।

\item Table-এর মাধ্যমে Data সুশৃঙ্খলভাবে রাখা যায়।

\item Data দ্রুত Search ও Process করা যায়।

\item একাধিক ব্যবহারকারী Database ব্যবহার করতে পারে।

\item Data-এর নিরাপত্তা নিয়ন্ত্রণের ব্যবস্থা রয়েছে।

\item Backup ও Recovery-এর ব্যবস্থা রয়েছে।

\item SQL ব্যবহার করে Database পরিচালনা করা যায়।

\item বৃহৎ প্রতিষ্ঠানের Database System-এ ব্যবহার করা যায়।

\end{compacttwocolpoints}

%================================================
% Uses
%================================================
\booksubheading{Oracle-এর ব্যবহার}

Oracle বিভিন্ন ধরনের Data-intensive System-এ ব্যবহার করা হয়।
যেমন—

\begin{compacttwocolpoints}

\item Banking System

\item Educational Database

\item Business Organization

\item Employee Management System

\item Inventory Management

\item Financial System

\item বড় প্রতিষ্ঠানের Information System

\end{compacttwocolpoints}

%================================================
% Oracle Programming Language?
%================================================
\booksubheading{Oracle কি Programming Language?}

\begin{rememberbox}

\textbf{\color{deepOrange}গুরুত্বপূর্ণ}

\textbf{না। Oracle নিজে Java বা Python-এর মতো General Purpose
Programming Language নয়।}

Oracle মূলত একটি—

\begin{center}

\textbf{
RDBMS / Database Platform
}

\end{center}

Oracle-এর সঙ্গে Database Programming ও Query-এর জন্য
\textbf{SQL ও PL/SQL} ব্যবহার করা হয়।

\end{rememberbox}

%================================================
% Java Python vs Oracle
%================================================
\booksubheading{Java/Python ও Oracle-এর পার্থক্য}

\begin{tabularx}{\linewidth}{|Y|Y|}
\hline

\rowcolor{softBlue}
\multicolumn{1}{|c|}{\textbf{Java / Python}}
&
\multicolumn{1}{c|}{\textbf{Oracle}}
\\
\hline

Programming Language
&
Database Management System
\\
\hline

Program ও Application তৈরিতে ব্যবহৃত হয়
&
Database তৈরি ও পরিচালনায় ব্যবহৃত হয়
\\
\hline

Programming Logic প্রকাশ করা হয়
&
Data সংরক্ষণ ও ব্যবস্থাপনায় প্রধান গুরুত্ব
\\
\hline

General Purpose Programming-এ ব্যবহার করা যায়
&
Database-oriented কাজে ব্যবহৃত হয়
\\
\hline

\end{tabularx}

%================================================
% Relationship
%================================================
\booksubheading{Oracle, SQL ও PL/SQL-এর সম্পর্ক}

Oracle, SQL ও PL/SQL-এর সম্পর্ক সহজভাবে—

\begin{center}
\begin{adjustbox}{max width=0.98\linewidth}
\fcolorbox{darkBlue}{softBlue}{
    \parbox[c][8mm][c]{0.22\linewidth}{
        \centering\textbf{Oracle Database}
    }
}
\quad
\(\longrightarrow\)
\quad
\fcolorbox{mainTeal}{softGreen}{
    \parbox[c][8mm][c]{0.24\linewidth}{
        \centering\textbf{SQL দিয়ে\\Data পরিচালনা}
    }
}
\quad
\(\longrightarrow\)
\quad
\fcolorbox{darkBlue}{softBlue}{
    \parbox[c][8mm][c]{0.26\linewidth}{
        \centering\textbf{PL/SQL দিয়ে\\Programming Logic}
    }
}
\end{adjustbox}
\end{center}

অর্থাৎ Oracle হলো Database Environment, SQL হলো Database-এর
Data নিয়ে কাজ করার ভাষা এবং PL/SQL হলো SQL-এর সঙ্গে Procedural
Programming সুবিধা যুক্ত ব্যবস্থা।

%================================================
% Final Reminder
%================================================
\begin{rememberbox}

\textbf{\color{deepOrange}মনে রাখবে}

\begin{itemize}

\item Oracle একটি \textbf{RDBMS}।

\item RDBMS-এর পূর্ণরূপ
\textbf{Relational Database Management System}।

\item Oracle Java বা Python-এর মতো
\textbf{General Purpose Programming Language নয়}।

\item Oracle Database-এ \textbf{SQL} ব্যবহার করা হয়।

\item SQL-এর পূর্ণরূপ
\textbf{Structured Query Language}।

\item PL/SQL-এর পূর্ণরূপ
\textbf{Procedural Language/SQL}।

\item PL/SQL-এ Variable, Condition, Loop, Procedure ও Function
ব্যবহার করা যায়।

\item SQL মূলত Data Query ও Database Operation-এর জন্য ব্যবহৃত হয়।

\item Oracle বড় Database System পরিচালনায় ব্যবহৃত হয়।
\item পরীক্ষায়
\textbf{Oracle = Programming Language}
লিখবে না।

\item সঠিক ধারণা—
\textbf{Oracle = RDBMS / Database Platform}।

\end{itemize}

\end{rememberbox}

\small

%================================================
% 5.9 Python
%================================================
